\documentclass[11pt]{article}
\usepackage[margin=1in]{geometry}
\usepackage{amsmath,amssymb,amsthm}
\usepackage{array,longtable,booktabs}
\usepackage[hidelinks]{hyperref}

\newtheorem{theorem}{Theorem}
\newtheorem{proposition}[theorem]{Proposition}
\newtheorem{corollary}[theorem]{Corollary}
\newtheorem{example}[theorem]{Example}
\newcommand{\Ann}{\operatorname{Ann}}
\newcommand{\diag}{\operatorname{diag}}

\title{P6-T06: A Scoped Gate-B Uniqueness, Robustness, and Admissibility Obstruction}
\author{The AEther-Flow Research Project}
\date{2026-07-26}

\begin{document}
\maketitle

\begin{center}
\fbox{\parbox{0.92\textwidth}{
\textbf{Status and authority.}
This is a task-local \emph{draft/control} Refuter result. It proves exact
comparison-family and transition countermodels and records a scoped
obstruction for the unchanged one-ray, unnormalized-density, dimensionless-
phase route. The comparison forms are not physical metrics. This result is
not a canonical-ontology edit, source-law adoption or rejection, causal-cone
identification, matter model, Gate Chair verdict, global no-go theorem,
publication, or completed derivation.
}}
\end{center}

\section{Question and fixed input}

P6-T03 leaves as its complete unconditional geometric input
\[
  \mathcal D_1(S,V)=\bigl([V]_+,\Ann(V),\preceq_V\bigr).
\]
P6-T04 enlarges this to
\[
  \mathcal D_{\rm cal}
  =\bigl(\mathcal D_1,\mu,\{\theta_A\}_{A\in\mathcal A}\bigr),
\]
where \(\mu\) is an independent positive source density and the
\(\theta_A\) are dimensionless scalar phases. P6-T05 proves that this package
does not select signature or nondegeneracy and that, on an admitted maximally
symmetric datum, source naturality forces every symmetric bilinear output to
vanish.

P6-T06 tests the remaining Gate-B burdens:
\begin{enumerate}
\item uniqueness up to a declared and observable-preserving gauge relation;
\item source naturality, transition compatibility, inverse laws, and cocycles;
\item robustness under finite variation and coarse-graining; and
\item operational admissibility for cones, clocks, rods, and universal
coupling.
\end{enumerate}
No target atlas, target metric, GR null cone, detector semantics, matter
response, \(G\), \(c\), or benchmark fit is admitted. The historical scoped
\(g_{\rm eff}\) source-extension object is not a premise and is unchanged.

\section{Typed assumption audit}

\begin{longtable}{>{\raggedright\arraybackslash}p{0.10\textwidth}
                  >{\raggedright\arraybackslash}p{0.38\textwidth}
                  >{\raggedright\arraybackslash}p{0.41\textwidth}}
\toprule
ID & Test assumption & Status and consequence\\
\midrule
\endhead
A1 & \(S\) is smooth and four dimensional; \(V\) is nonzero and positively
oriented. & Inherited candidate scope, not a physical spacetime or clock
conclusion.\\
A2 & Only \(\mathcal D_{\rm cal}\) is unconditional. & No coframe, bilinear
form, connection, target chart, operational transport, or universality law is
hidden in the input.\\
A3 & \(\mu\) is independent source density data. & It is not identified with
the volume density of a comparison form.\\
A4 & Scalar phases may be constant on the exact comparison datum. & They then
provide no discriminator while remaining admitted data.\\
A5 & Comparison forms are mathematical controls on one fixed source basis. &
They are neither source geometry nor target geometry.\\
A6 & A gauge quotient must make stated observables representative-independent.
& Merely naming a source automorphism ``gauge'' is not an operational
transport theorem.\\
A7 & P6-T03 through P6-T05 are fixed inputs. & Their scope is not enlarged.\\
A8 & Generated artifacts, validators, roles, and route authority are not
scientific premises. & Process success cannot establish geometry.\\
\bottomrule
\end{longtable}

\section{The same-source comparison family}

Fix a source tangent basis \(e_0,e_1,e_2,e_3\), the oriented ray
\(L=[e_0]_+\), the density
\(\mu_0=|e^0\wedge e^1\wedge e^2\wedge e^3|\), and constant phases. For
\(a\in\mathbb R\) and \(\varepsilon>0\), define
\[
 q_{a,\varepsilon}
 =\diag\!\left(-1,e^{2a},e^{-2a},\varepsilon^2\right).
\]
Every member has Lorentzian inertia and
\[
 \det q_{a,\varepsilon}=-\varepsilon^2.
\]
The admitted source carriers are independent of \(a\) and
\(\varepsilon\).

\begin{theorem}[Cone and transverse-comparison nonselection]
\label{thm:cone}
The family \(q_{a,\varepsilon}\) changes both the quadratic classification of
a fixed source covector and transverse comparisons while the admitted source
datum stays fixed.
\end{theorem}
\begin{proof}
The inverse matrix is
\[
 q_{a,\varepsilon}^{-1}
 =\diag\!\left(-1,e^{-2a},e^{2a},\varepsilon^{-2}\right).
\]
For the fixed source covector \(\xi=e^0+e^1\),
\[
 q_{a,\varepsilon}^{-1}(\xi,\xi)=-1+e^{-2a}.
\]
It is positive for \(a<0\), zero for \(a=0\), and negative for \(a>0\).
Meanwhile
\[
 q_{a,\varepsilon}(e_1,e_1)=e^{2a},\qquad
 q_{a,\varepsilon}(e_2,e_2)=e^{-2a}.
\]
Thus cone classification and transverse comparisons both vary at fixed
source data.
\end{proof}

\begin{corollary}[A gauge declaration does not close the burden]
\label{cor:gauge}
Let
\[
 B_a=\diag(1,e^a,e^{-a},1).
\]
Then \(\det B_a=1\), \(B_a\) preserves the oriented ray and density, and
\(q_{a,\varepsilon}=B_a^Tq_{0,\varepsilon}B_a\). Treating these
source-data automorphisms as gauge does not by itself define a physical cone.
\end{corollary}
\begin{proof}
The matrix identities are direct. If forms are quotient identified, a
representative-independent cone statement also needs a derived identification
of operational covectors between representatives. The unchanged source
package derives neither that transport nor an invariant operational
readout. Indeed Theorem~\ref{thm:cone} shows that the same named source
covector has different classifications across representatives. Co-transforming
the covector restores the elementary tensor identity, but that conditional
move is precisely the missing source-derived transport law.
\end{proof}

\section{Finite-variation and conditioning stress}

\begin{proposition}[Nonzero response at fixed source input]
\[
 \frac{\partial q_{a,\varepsilon}}{\partial a}
 =\diag\!\left(0,2e^{2a},-2e^{-2a},0\right)\ne0.
\]
No variation of \(\mathcal D_{\rm cal}\) accounts for this response.
\end{proposition}

\begin{proposition}[No uniform nondegeneracy certificate]
As \(\varepsilon\to0^+\), the inverse component
\((q_{a,\varepsilon}^{-1})^{33}=\varepsilon^{-2}\) diverges and the limiting
form is degenerate. The signed family
\[
 r_t=\diag(-1,1,1,t)
\]
is Lorentzian for \(t>0\), degenerate at \(t=0\), and split-signature for
\(t<0\). Current source data derive no positive lower bound or exclusion of
this crossing.
\end{proposition}

P5 source-side error and semigroup certificates do not change this
conclusion. A transfer theorem would require a geometry-valued map
\(\mathfrak G\), a comparison distance \(d_q\), and an estimate such as
\[
 d_q\!\left(\mathfrak G(C_\ell D),\mathfrak G(D)\right)
 \le L\,d_D(C_\ell D,D).
\]
None of \(\mathfrak G,d_q,L\) is derived. Source-side stability therefore
does not certify geometry-side stability.

\section{Covariance, inverse, and cocycle controls}

P6-T05 already proves that the determinant-one diagonal automorphisms of an
admitted maximally symmetric datum force the only natural symmetric bilinear
form to be zero. P6-T06 does not re-prove or broaden that theorem; it uses it
as the cumulative naturality obstruction.

The following finite controls separate local signature from global
covariance. With
\(\eta=\diag(-1,1,1,1)\), the transition \(2I\) fails because
\[
 (2I)^T\eta(2I)=4\eta\ne\eta.
\]
Even when every transition is Lorentz-valued, inverse and cocycle laws remain
independent. The pair
\(\Lambda_{VU}=I,\Lambda_{UV}=-I\) violates
\(\Lambda_{UV}=\Lambda_{VU}^{-1}\). On a triple overlap,
\[
 \Lambda_{VU}=I,\qquad \Lambda_{WV}=I,\qquad
 \Lambda_{WU}=-I
\]
has Lorentz-valued entries but violates
\(\Lambda_{WV}\Lambda_{VU}=\Lambda_{WU}\).
These are conditional consistency countermodels, not claims that the current
source package derives any transition maps.

\section{Physical-admissibility audit}

\begin{longtable}{>{\raggedright\arraybackslash}p{0.27\textwidth}
                  >{\raggedright\arraybackslash}p{0.48\textwidth}
                  >{\raggedright\arraybackslash}p{0.15\textwidth}}
\toprule
Criterion & Exact audit & Disposition\\
\midrule
\endhead
Unique cone up to gauge & Theorem~\ref{thm:cone} and
Corollary~\ref{cor:gauge}; no operational transport or quotient-observable
theorem. & Not met\\
Source covariance & P6-T05 naturality obstruction; transition, inverse, and
cocycle controls above. & Not met\\
Finite-variation robustness & Inverse conditioning grows as
\(\varepsilon^{-2}\); degeneracy and signature crossing are not excluded. &
Not met\\
Characteristic cone & P6-T02 gives \(\Ann(V)\), a linear hyperplane, not a
nondegenerate Lorentzian quadratic null cone. & Not met\\
Clock calibration & P6-T04 leaves positive rescaling \(V_f=fV\) and no
absolute normalization. & Not met\\
Rod calibration & Transverse comparisons vary as \(e^{2a}\) and \(e^{-2a}\)
at fixed source input. & Not met\\
Coarse-grained stability & No geometry map, output distance, or transfer
estimate is derived. & Not met\\
Universal matter coupling & No detector semantics, matter action,
universality theorem, or exclusion of bimetric, multicone, and nonuniversal
comparison branches is derived. & Not met\\
\bottomrule
\end{longtable}

The last row does not adopt any bimetric, multicone, or nonuniversal physical
sector. Those are only unexcluded comparison branches.

\section{Precise obstruction and route consequence}

\begin{theorem}[Scoped Gate-B obstruction]
\label{thm:obstruction}
For the exact cumulative source package tested, Gate B is not closed. The
package does not determine a unique cone modulo a derived observable-
preserving gauge relation, a nonzero natural bilinear form, globally
compatible transitions, uniform nondegeneracy, geometry-level
coarse-graining stability, characteristic-cone identification, calibrated
clocks and rods, or universal operational coupling.
\end{theorem}
\begin{proof}
Cone and rod nonselection follow from Theorem~\ref{thm:cone}. The gauge burden
follows from Corollary~\ref{cor:gauge}. Naturality is obstructed by P6-T05.
The transition controls show that local signature does not supply
Lorentz-valued transitions, inverses, or cocycles. The
\(\varepsilon\)-family gives the missing uniform bound. The
physical-admissibility table records the remaining exact non-implications.
\end{proof}

The obstruction identifier is
\[
\texttt{OBST-P6T06-GATE-B-UNIQUENESS-ROBUSTNESS-ADMISSIBILITY-001}.
\]
Its Refuter classification is \texttt{scoped\_obstruction}. The exact
unchanged route from \(\mathcal D_{\rm cal}\) to an unscoped physical geometry
is locally frozen. This is not a global no-go theorem. A materially richer
source-side law could supply a discriminator, operational transport, uniform
robustness, a geometry-level coarse-graining map, and a universality theorem.
Current status is
\texttt{blocked\_adoption\_open\_continuation}.

\section{Distance-to-GR disposition}

P6-T06 supplies new countermodels and a precise obstruction, so it is not a
route orbit. Gate B remains unpassed. The historical scoped \(g_{\rm eff}\)
object is unchanged. Physical causal geometry, unscoped \(g_{\rm eff}\),
matter coupling, Einstein equations, benchmark promotion, proof authority,
publication, push, and completed derivation remain blocked.

\section{Project source note}

The fixed project basis is the registered P6-T03, P6-T04, and P6-T05 theorem
packets, canonical target-side geometry dictionary, historical scoped
source-extension \(g_{\rm eff}\) candidate and Gate Chair record, v21
implementation plan, Distance-to-GR burden map, and handoff-0883. Exact
source hashes are recorded in the accompanying machine-readable
specification. No external physical datum is used.

\end{document}
